% !TeX root = ../../surface-lagrangian-torus-fibration.tex
% 2026-11-21 / v1.0: Source を含む法線方向の係数一と正則深度障壁．
\section{Source を含む法線方向の係数一の評価}
\label{ke:sharp-normal-section}

本節では actual KE 計量に対する一方向の評価を示す．
固定した残余 compact 上で正則な normal frame を用い，
source を除かずに最大値原理を適用する．外側の固定 shell では
有限な上界だけを用いる．その値の大きさは，最後に選ぶ
global depth の下側 cut に吸収する．従って成分 KE 計量の
sharp normal asymptotic は本節の追加入力ではない．
全ての評価は平均化前の actual $\omega_t$ に対するものであり，
非正則な Moser map を介して正則性を移す議論は用いない．

\subsection{規格化と一変数の障壁}

$\operatorname{Ric}(\omega_t)=-\omega_t$，
$L=-\log|t|^2$，$R=\log L$ とする．
$\omega_t=\sqrt{-1}g_{a\bar b}\,dz^a\wedge d\bar z^b$ に対し，
\[
 \Delta_{\omega_t}=g^{a\bar b}\partial_a\partial_{\bar b},
 \qquad |df|_{g_t}^2=2|\partial f|_{\omega_t}^2
\]
を用いる．非零正則関数 $x$ が $0<|x|<1$ を満たす領域で
\[
 s=-\log|x|^2>0,\qquad
 h_{\mathrm{cusp}}=dd^c(-2\log s)
 =\frac{2\sqrt{-1}\,dx\wedge d\bar x}{|x|^2s^2},
 \qquad U=\operatorname{tr}_{\omega_t}(x^*h_{\mathrm{cusp}})
       =|d\log s|_{g_t}^2
       =\frac{2|\partial s|_{\omega_t}^2}{s^2}
\]
と置く．$h_{\mathrm{cusp}}$ は Ricci curvature $-1$ の
punctured-disc 計量である．Rank-one Chern--Lu の計算\cite[Theorem~2およびその証明]{YauSchwarz}は
\begin{equation}\label{ke:eq:sharp-normal-logistic}
 \Delta_{\omega_t}\log U\ge U-1,
 \qquad \Delta_{\omega_t}U\ge U(U-1)
\end{equation}
を与える．$dx\ne0$ の場合は $U>0$ なので通常の微分不等式である．
この係数は本稿の $d^c=\frac{\sqrt{-1}}2(\bar\partial-\partial)$ と
$\operatorname{Ric}(\omega_t)=-\omega_t$ の規格化による．
実際，$\partial s=-dx/x$ は正則 one-form なので，
Bochner 計算と Cauchy--Schwarz から
$\Delta\log|\partial s|_{\omega_t}^2\ge-1$ となる．
また $\Delta s=0$ より
$\Delta(-2\log s)=2|\partial s|_{\omega_t}^2/s^2=U$ である．
両式を加えると第一不等式，さらに
$\Delta U=U\Delta\log U+|\partial U|^2/U$ から第二不等式を得る．

\begin{lemma}[二つの端の粗い上界を吸収する障壁]
\label{ke:normal-harmonic-depth-barrier}
比較的 compact な領域 $\mathcal A$ で上記の $x,s,U$ を考える．
境界を $s=n$，$s=T$，および残余境界 $\Gamma$ に分け，
$0<n<T$ とする．境界が corners を持つことを許す．
ある $\varepsilon,A,C>0$ に対して
\[
 U|_{s=n}\le A,\qquad U|_{s=T}\le C,
 \qquad U|_\Gamma\le1+\varepsilon
\]
ならば，全 $\mathcal A$ で
\begin{equation}\label{ke:eq:normal-harmonic-depth-barrier}
 U\le1+\varepsilon+A\frac ns+C\frac sT
\end{equation}
が成り立つ．
\end{lemma}

\begin{proof}
$s$ は pluriharmonic なので $\Delta s=0$ である．
\[
 B(s)=1+\varepsilon+A\frac ns+C\frac sT
\]
と置くと，$B\ge U$ が境界で成立し，
\[
 \Delta B=B''(s)|\partial s|_{\omega_t}^2
          =\frac{An}{s}U
\]
である．$U-B$ が正の最大値を内点で取ると，その点で
\[
 0\ge\Delta(U-B)
 \ge U\left(U-1-\frac{An}{s}\right)
 >U\left(\varepsilon+C\frac sT\right)>0
\]
となり矛盾する．境界の corners は内点の最大値原理に影響しない．
\end{proof}

\subsection{Source を除かない正則な比較領域}

二重線 $D=D_{ik}$ と，節\ref{ke:neck-corner-transition-section}で
選んだ許容射影を固定する．三重点の集合を $Q\subset D$ とする．
一つの $Q$ の点を無限遠とする affine 座標を取り，
$D\setminus Q$ 上の二つの normal line bundles を自明化する．
これは定理\ref{ke:projected-source-neck}の証明，および
節\ref{ke:neck-corner-exhaustion-section}の global frame と同じ選択である．
射影 Morse 変換により
\begin{equation}\label{ke:eq:sharp-normal-global-chart}
 x'_g y'_g=t h_t(p),\qquad p\in D\setminus Q
\end{equation}
と書ける．Frame は $t$ に依存せず，Morse 変換は固定残余 compact
上の固定幅で恒等写像から全固定階で $O(t)$ ずれる．

$b_j\to\infty$ を取り，三重点近傍の固定座標
$Z_{a,0}$ を使って，境界が
$-\log|Z_{a,0}|^2=b_j$ となる残余 compact $P_j$ を取る．
有限個の初期段階を除き，全 source は $P_j$ の内部にある．
実際の construction に使うより大きい残余領域とは区別し，
本節の $P_j$ は最大値原理の補助領域とする．
少し広い $P_j^+$ も先に固定する．

固定した $0<\kappa<1/4$ は，補題\ref{ke:source-arm-scale-buffer}の
$\theta$ より十分小さく取る．$n_j$ を後に大きく選び，
\begin{equation}\label{ke:eq:sharp-normal-domain}
 \mathcal A^x_{j,t}
 =\{p\in P_j,\ n_j\le s=-\log|x'_g|^2\le\kappa L\},
 \qquad y'_g=\frac{t h_t(p)}{x'_g}
\end{equation}
とする．$n_j$ が十分大きく，その後 $|t|$ が十分小さければ，
これは全角度を含む annulus と $P_j$ の積である．
実際，$|x'_g|\le e^{-n_j/2}$ であり，
\[
 |y'_g|\le\sup_{P_j}|h_t|\,e^{-(1-\kappa)L/2}\longrightarrow0.
\]
従って全領域が固定幅の normal chart 内に入り，
$y'_g$ の別の境界は現れない．$h_t(p)=0$ でも $x'_g\ne0$ なので，
$(x'_g,p)$ は正則座標であり，source の零点を $P_j$ から除く必要はない．
対称な $y$-arm の領域 $\mathcal A^y_{j,t}$ も同様に定める．

\begin{lemma}[比較領域の三種類の境界評価]
\label{ke:sharp-normal-boundary-inputs}
任意の $\varepsilon_j\downarrow0$ に対し，$b_j$ を十分大きく，
次に $n_j$ を十分大きく選べる．その後に有限な定数 $A_j,C_j$ があり，
$|t|$ が十分小さければ $\mathcal A^x_{j,t}$ 上の $U$ は
\begin{equation}\label{ke:eq:sharp-normal-boundaries}
 U|_{s=n_j}\le A_j,\qquad
 U|_{s=\kappa L}\le C_j,\qquad
 U|_{p\in\partial P_j}\le1+\varepsilon_j
\end{equation}
を満たす．有限個の二重線と両向きの arm に対して同時に選べる．
\end{lemma}

\begin{proof}
\emph{外側の固定 shell.}
$P_j$ と $n_j$ を固定すると，$s=n_j$ は $t=0$ の成分の smooth locus
の compact 部分へ収束する．Source の基底点でも $x'_g\ne0$ なので同じである．
固定 gauge における KE 計量の smooth convergence と
Morse 座標の $O(t)$ 収束から有限な $A_j$ を得る．
$A_j$ が $j$ 又は $n_j$ とともに有界である必要はない．

\emph{残余境界.}
三重点 end の一つで，局所 normal frame と global frame の比を
非零正則 base-only function $a_j(p)$ と書く．
補題\ref{ke:arm-coordinate-ratio}と命題\ref{ke:arm-log-derivatives}により
\[
 x'_g=a_j(p)C_D(p,t)\zeta_i\exp(e_{ki}),
 \qquad
 s=\ell_i^{\rm tr}+f_{j,t}(p)-2\Re e_{ki},
 \quad f_{j,t}=-\log|a_jC_D|^2.
\]
$\partial P_j$ の固定した少し広い annuli 上で，
$|f_{j,t}|+|df_{j,t}|_{g_t}\le D_j$ となる有限な $D_j$ を取れる．
ここで関数自体の微分は固定 ambient coordinates で有界であり，
定理\ref{ke:fixed-base-neck-moser}の証明の射影下界によって
actual covector norm も一様に抑えられる．この定数は $n_j,L$ に依存しない．

まず $n_j\gg b_j+D_j$ とする．固定 end annulus の frame 定数も $n_j$ に吸収すると，
元の射影座標で $|x|+|y|\ll|p|$ となる．また $|t|\ll|p|^3$ は最後の $t$ 閾値で
満たせるため，座標比の補題を arm の深度比較に先立って適用できる．$\partial P_j$ では
$\ell_k^{\rm tr}=b_j+o_t(1)+O(e^{-c b_j})$ であり，
$n_j\le s\le\kappa L$ から残りの二深度は
$\ell_i^{\rm tr},\ell_k^{\rm tr}$ の関係を固定比で分離する．
より正確には，端の残余添字を $k$ とし，他方の normal 添字を $i'$ とすると
\[
 \ell_i^{\rm tr},\ell_{i'}^{\rm tr}
 \ge(1+\sigma)\ell_k^{\rm tr}
\]
がある固定 $\sigma>0$ について成立する．第二の不等式には
$\ell_{i'}^{\rm tr}=L-\ell_i^{\rm tr}-\ell_k^{\rm tr}$ と，
$t$ を最後に小さくすることを用いる．
従って上記の座標比と全固定階の評価をこの領域全体で使える．

関連 KE 主稿\cite{KELocalCurrent}の
\texttt{simplex:lem:rank-one-schwarz} と
\texttt{coarse:thm:full-triple-tensor} は
\[
 |d\log\ell_i^{\rm tr}|_{g_t}\le1+O(b_j^{-1})
\]
を与える．また同主稿の complete-simplex と cusp reference metric の
一様比較により，命題\ref{ke:arm-log-derivatives}の誤差を
$C^1(g_t)$ でも $O(e^{-c b_j})$ と評価できる．
従って
\[
 |d\log s|_{g_t}
 \le \left(1+\frac{D_j+Ce^{-c b_j}}{n_j}\right)
             (1+Cb_j^{-1})
       +\frac{D_j+Ce^{-c b_j}}{n_j}.
\]
$b_j$，次に $n_j$ を選べば，この二乗を $1+\varepsilon_j$ 以下にできる．
これは正則な座標間の比較であり，三つの arm の smooth 補間を用いていない．

\emph{内側の moving shell.}
Source の固定した小円板の外では，
定理\ref{ke:fixed-base-neck-moser}の annulus--disc 二側比較から，
$s=\kappa L$ 上で $U\le C_j$ を得る．
Source の内側では $z=h_t(p)$ を局所座標に取り，
同じ normal frame において $x'_g y'_g=tz$ を使う．
必要な固定 rescaling のため深度に有界な定数が加わることを許す．
補題\ref{ke:source-arm-scale-buffer}の positive function は，
$S=s+O_j(1)=\kappa L+O_j(1)$ で
\[
 u_t=F_t+2R+C_*
     =2\log(L/S)+O(1)=2\log(1/\kappa)+O_j(1).
\]
従って同補題は，中心を $s=\kappa L$ に取った全ての点で，
ある $\rho_j>0$ を半径とする actual 球が，正しい arm の
$0<|x'_g|<1$ の領域に含まれることを与える．
これは min-depth の下界だけでなく，source の中央 barrier を横切らない
orientation の保持を含む．
この球上で正則写像 $x'_g$ への rank-one localized Chern--Lu を適用する．
ここで必要なのは，同 KE 主稿 \texttt{coarse:lem:localized-chern-lu} の
cutoff 証明を \eqref{ke:eq:sharp-normal-logistic} に適用することであり，
rank-one pullback を非退化な K\"ahler form とみなすことではない．
半径 $\rho_j$ の球内に支持され，半径 $\rho_j/2$ で $1$ となる
actual-distance cutoff $\chi$ を取る．Ricci 下界と Laplacian comparison により
$|\partial\chi|^2\le C\rho_j^{-2}$，
$\Delta\chi\ge-C(\rho_j^{-2}+\rho_j^{-1})$ を barrier の意味で満たせる．
$\chi^2U$ の正の最大点で
\[
 0\ge\Delta\log(\chi^2U)
 \ge U-1+2\frac{\Delta\chi}{\chi}
                  -2\frac{|\partial\chi|^2}{\chi^2}
\]
に $\chi^2$ を掛けると，この最大点で
\[
 \chi^2U\le\chi^2-2\chi\Delta\chi+2|\partial\chi|^2
       \le C(1+\rho_j^{-2})
\]
となるので $U\le C(1+\rho_j^{-2})$ を得る．
Cut locus では同じ support-function argument
を用いる．
Source は有限個なので上界を一つの $C_j$ にまとめられる．
\end{proof}

\begin{proposition}[Source を含む sharp scalar normal bound]
\label{ke:sharp-scalar-normal-barrier}
上記の選択で，$|t|$ が十分小さければ
\begin{equation}\label{ke:eq:sharp-scalar-normal-barrier}
 |d\log s|_{g_t}^2
 \le1+\varepsilon_j+A_j\frac{n_j}{s}
                  +C_j\frac{s}{\kappa L}
 \quad\text{on }\mathcal A^x_{j,t}.
\end{equation}
同じ評価は $y$-arm にも成立する．
\end{proposition}

\begin{proof}
補題\ref{ke:sharp-normal-boundary-inputs}を
補題\ref{ke:normal-harmonic-depth-barrier}に適用する．
\end{proof}

\subsection{Global physical depth と active cut}

Global Hermitian norm で定めた
$\ell_i=-\log\|L_i\|^2$ と，正則 frame の $s$ を区別する．
下側 cut $\mathfrak m_j$ は旧 closed-corner の cut $6b_j$ と別に選ぶ．
上側 cut を $M_{j,t}=L/K_j$ とする．

\begin{lemma}[Active band 上の physical depth と正則深度]
\label{ke:sharp-normal-physical-unit}
各固定 $j$ について有限な $D'_j$ がある．
$K_j>4$，$\mathfrak m_j$ が十分大きいとき，
対応する pair neighborhood 内の
\[
 \mathfrak m_j\le\ell_i\le L/K_j
\]
では，$t$ を十分小さくすれば
\begin{equation}\label{ke:eq:sharp-normal-physical-unit}
 |\ell_i-s|+|d(\ell_i-s)|_{g_t}\le D'_j.
\end{equation}
定数は $t,\mathfrak m_j,K_j$ に依存せず，固定した残余 compact と
frame に依存してよい．
\end{lemma}

\begin{proof}
固定残余 compact 上では，global bundle frame の元の法線座標 $x_g$ と
平面 section の局所表示の間は，非零な固定 base factor と smooth
Hermitian norm の項だけ異なる．これらの対数と微分は固定領域で有界である．
射影 Morse 変換について $x'_g=x_g+O(t)$ が全固定階で成立する．
Active upper cut より $|x_g|\ge c_j|t|^{1/K_j}$ なので
\[
 \left|\frac{x'_g-x_g}{x_g}\right|
       \le C_j|t|^{1-1/K_j},\qquad
 \left|d\log\frac{x'_g}{x_g}\right|_{g_{\mathrm{FS}}}
       \le C_j|t|^{1-2/K_j}=o_t(1).
\]
固定 residual 部分の射影下界
$g_t\ge c_jg_{\mathrm{FS}}|_{X_t}$ は，source-free 部分では
定理\ref{ke:fixed-base-neck-moser}の証明，source 近傍では
\eqref{ke:eq:source-global-positive-u}から得る．
従って上の誤差の actual covector norm も有界であり，
Hermitian frame 項と合わせて所要の $D'_j$ を得る．
$O(t)$ という絶対誤差を深度にそのまま代入せず，
active upper cut による相対誤差の評価をここで使用した．
\end{proof}

\begin{lemma}[深い三重点域の外側の active 点の配置]
\label{ke:sharp-active-pair-cover}
$\mathcal T_{j,t}$ を，三つ以上の global physical depths が
$b_j/2$ 以上である領域とする．$j$ が十分大きく，$K_j>2d$ とする．
$\mathfrak m_j\gg b_j$ を固定して，最後に $t$ を十分小さくすると，
$\mathcal T_{j,t}$ の外で
$\mathfrak m_j<\ell_i<L/K_j$ を満たす点は，ある $k\ne i$ に対する
上記 $D_{ik}$ の pair neighborhood に入り，その射影は $P_j$ にある．
その点の正則深度は
\[
 s\ge\mathfrak m_j-D'_j,\qquad s\le L/K_j+D'_j.
\]
\end{lemma}

\begin{proof}
族の方程式と固定 Hermitian norm から
\[
 \sum_{a=1}^d\ell_a=L-\log\|F_d\|^2\ge L-C
\]
である．従ってある $k$ で $\ell_k\ge(L-C)/d>L/K_j$ となる．
$\ell_i<L/K_j$ なので $k\ne i$ である．
一方 $\mathcal T_{j,t}$ の外では他の全深度が $b_j/2$ 未満であり，
点は $D_{ik}$ の近傍にあるが，三重点の小さい端には入らない．

この最後の包含を定量的に確認する．有限個の fixed projective charts
上で二つの normal 座標の大きさは
$O(e^{-\mathfrak m_j/2})$ と $O(e^{-cL})$ である．
第三平面の section の，点における値と射影点における値の差は
この和の定数倍以下である．第三深度が $b_j/2$ 未満なので，
その値は固定 frame で $c e^{-b_j/4}$ 以上である．
$\mathfrak m_j$ を $b_j$ より十分大きく取れば，相対誤差は小さい．
したがって射影点の end depth は $b_j/2+O(1)$ 以下であり，
cut を $b_j$ に置いた $P_j$ の内側にある．
固定 projective charts と end 座標の単位による誤差は有界で，
$b_j/2$ と $b_j$ の余裕に吸収できる．
これにより fixed pair chart とその projection が定義される．
正則深度の二つの評価は
補題\ref{ke:sharp-normal-physical-unit}による．
\end{proof}

\begin{theorem}[Global one-active covector の係数一の上界]
\label{ke:sharp-global-one-active}
任意の $\varepsilon_j\downarrow0$ に対し，固定段階の選択を
\[
 b_j\ ,\ P_j\ ,\ n_j\ ,\ A_j,C_j,D'_j\ ,\
 \mathfrak m_j\ ,\ K_j\ ,\ t_j
\]
の順に行える．ここで $\mathfrak m_j\to\infty$，$K_j\to\infty$ とし，
例えば
\begin{equation}\label{ke:eq:sharp-normal-parameter-order}
 \frac{A_jn_j+D'_j+n_j}{\mathfrak m_j}\le\varepsilon_j,
 \qquad \frac{C_j}{\kappa K_j}\le\varepsilon_j,
 \qquad K_j>\max\{2d,4,4/\kappa\}
\end{equation}
を課す．さらに $\mathfrak m_j$ を固定 stage の component core の
non-dominant depth 上界より大きくしてよい．
$|t|\le t_j$ とすると，全ての点と添字に対して
\begin{equation}\label{ke:eq:sharp-global-one-active}
 \mathfrak m_j<\ell_i<L/K_j
 \quad\Longrightarrow\quad
 |d\log\ell_i|_{g_t}\le1+C\varepsilon_j+C/b_j
\end{equation}
が成立する．定数 $C$ は固定 $d$ と族に依存し，$j,t$ に依存しない．
\end{theorem}

\begin{proof}
$\mathcal T_{j,t}$ の外では
補題\ref{ke:sharp-active-pair-cover}により，対応する pair chart を取る．
$t$ を最後に小さくすれば
\[
 n_j\le s\le2L/K_j<\kappa L,
 \qquad s\ge\mathfrak m_j/2.
\]
命題\ref{ke:sharp-scalar-normal-barrier}から
\[
 |d\log s|_{g_t}^2
 \le1+\varepsilon_j+\frac{2A_jn_j}{\mathfrak m_j}
                    +\frac{2C_j}{\kappa K_j}
 \le1+5\varepsilon_j.
\]
補題\ref{ke:sharp-normal-physical-unit}により
\[
 |d\log\ell_i|_{g_t}
 \le\left(1+\frac{D'_j}{\mathfrak m_j}\right)|d\log s|_{g_t}
       +\frac{D'_j}{\mathfrak m_j},
\]
なので所要の上界を得る．

$\mathcal T_{j,t}$ の中では，平面の一般位置性により，
$b_j$ が十分大きければ一つの固定三重点 chart に入る．
三つの正確な深度は $b_j/2-O(1)$ 以上であり，
physical depth は正確な深度と，固定 smooth unit の対数だけ異なる．
Complete-simplex の係数一の rank-one Schwarz 評価と
actual/model $O(b_j^{-1})$ 比較を用いる．
固定 smooth unit の微分は cusp reference metric の covector norm で
有界であり，同 reference metric と actual metric の一様比較により
その actual norm も有界である．Active depth で割れば誤差は
$O(\mathfrak m_j^{-1})$ となり，同じ結論を得る．
有限個の全 charts で定数と最後の $t$ 閾値を共通に取る．
\end{proof}

\begin{remark}[遅い対角選択と適用範囲]
\label{ke:sharp-normal-diagonal-remark}
Stage $j$ の全定数を先に選んだ後で，
\[
 R_t\ge j\bigl(1+\log\mathfrak m_j+\log K_j\bigr)
\]
を追加してよい．従って $j=j(t)\to\infty$ を十分遅く選べば
\[
 \frac{\log\mathfrak m_{j(t)}+\log K_{j(t)}}{R_t}\longrightarrow0,
 \qquad
 S_{j,t}=\log\frac{L/K_j}{\mathfrak m_j}\sim R_t.
\]
これは \eqref{ke:eq:sharp-global-one-active} を clipped logarithmic map
\[
 \Theta_{i,t}
 =\left[\frac{\log(\ell_i/\mathfrak m_j)}{S_{j,t}}\right]_0^1
\]
へ適用するための選択である．二つの active covectors の同時評価は
full-triple の dominant-product 比較で別に行う．
本節の係数一の評価は，source を含む moving neck に対する
full tensor comparison，TC，PBE を仮定しない．
全曲面の距離下界は，この scalar 評価と二方向 co-Gram 評価を用いて
全曲線に沿って $\Theta_t$ の長さを評価することで得られる．
Fibration の値との整合と距離上界はそれぞれ別の主張である．
\end{remark}
